\section{Scaling to \texorpdfstring{$N$}{N} Hops}
\label{sec:nhop}

The concrete theory fixes the path at three hops. Whether the same guarantees
hold for paths of arbitrary length is the obvious next question, and it is not
answerable by simply adding rules: each additional hop in the signed-message
encoding multiplies the adversary's options, and \mtamarin's backward search
does not terminate.

\subsection{Abstracting the Routing Layer}

The obstruction is the signed HTLC offer. Resolving where a forwarded message
came from requires \mtamarin\ to consider every party as a possible origin, and
with an unbounded number of hops that analysis recurses without bound. We
therefore replace the network-level message with an internal linear fact handed
directly from one hop to the next:

\begin{lstlisting}
rule Forward_HTLC:
    [ Route($Prev, $Me, ptrIn, y, amtOut + fee),
      !ChannelConnect(ptrOut, $Me, $Next),
      Free(ptrOut) ]
    --[ HTLCForwarded($Me, $Prev, $Next, ptrIn, ptrOut, y,
                      amtOut + fee, amtOut),
        FeeExpected($Me, ptrIn, ptrOut, y, fee),
        Neq(ptrIn, ptrOut), Neq($Prev, $Me), Neq($Me, $Next) ]->
    [ IntermediaryPending($Me, $Prev, $Next, ptrIn, ptrOut, y,
                          amtOut + fee, amtOut, fee),
      Route($Me, $Next, ptrOut, y, amtOut) ]
\end{lstlisting}

This single rule fires at every hop, for any $N$. Because the HTLC is no longer
a self-referential signed term, \mtamarin's origin analysis no longer recurses.
The price is that per-hop authentication is no longer machine-checked at this
layer; \Cref{sec:refinement} discharges that obligation separately.

\subsection{Amount Chaining and the Choice of Builtin}

The rule above is worth reading closely for one detail. The incoming amount is
matched against the pattern \code{amtOut + fee}, where \code{+} is the
associative-commutative operator of the \code{multiset} builtin. AC unification
lets this pattern \emph{split} the incoming amount nondeterministically: $\$Me$
keeps \code{fee} and locks \code{amtOut} downstream. The outgoing \code{Route}
fact then carries \code{amtOut}, which becomes the next hop's incoming amount.
Value therefore chains along a path of any length.

This is why the two theories use different builtins. The
\code{natural-numbers} builtin used in \code{multihop.spthy} provides
\code{\%+}, which can \emph{check} a sum --- \code{Eq(\%vIn, \%vOut \%+ \%fee)}
--- but cannot destructure one in a premise, because \code{\%+} is not
AC-unifiable. At three fixed hops, checking suffices, and
\code{natural-numbers} keeps the arithmetic lemmas cheap. At unbounded $N$,
splitting is required, so \code{multiset} is the only workable choice. The
selection is driven by what each builtin can unify, not by preference.

Five properties are new at this layer and could not be stated at all without
the chaining: \code{Forward\_Requires\_Incoming} (279 steps), which requires
every forward to be preceded by an offer or a predecessor forward \emph{with
agreeing amounts}; \code{Amount\_Strictly\_Decreases\_Per\_Hop};
\code{Fee\_Only\_On\_Successful\_Forward};
\code{Redeem\_Requires\_Receiver\_Release}; and the witness
\code{Multihop\_Payment\_With\_Fees\_Possible}.

\subsection{Soundness of the Abstraction}
\label{sec:refinement}

The abstraction transfers safety from the concrete model to the abstract one
only if the concrete model's authentication guarantees justify it. We make this
precise. Let $C$ be the concrete theory and $A$ the abstract one, and let $\pi$
project a concrete trace onto the shared alphabet $\Sigma$ of routing events:
\[
\small
\Sigma = \left\{
\begin{array}{l}
\mathit{HTLCOffered},\ \mathit{HTLCForwarded},\\
\mathit{HTLCFulfilled},\ \mathit{Redeemed},\ \mathit{Refunded}
\end{array}
\right\}.
\]

\begin{thrm}[Refinement]
\label{thm:refinement}
Under the per-hop authentication guarantees of \Cref{thm:causality}, for every
honest trace $t$ of $C$ there is a trace $t'$ of $A$ with $\pi(t) = \pi(t')$.
Consequently every $\Sigma$-property proved of $A$ holds of $C$ restricted to
honest executions.
\end{thrm}

The argument is by induction on the length of the concrete trace and is given
in full in the artifact. The essential step is that \Cref{thm:causality} lets a
concrete \code{Forward1\_HTLC} step be matched by an abstract
\code{Forward\_HTLC} step consuming the corresponding \code{Route} fact: the
signed message's authenticity is exactly what the linear fact assumes.

\Cref{thm:refinement} is a pen-and-paper result, not a machine-checked one, and
we flag it as the principal soundness caveat of the $N$-hop extension.

\medbfparagraph{Why the wormhole disappears}
One consequence deserves comment. The wormhole is not reachable in the abstract
model, and this is structural rather than accidental. In $C$ the preimage
travels as a network message that the adversary can relay out of band. In $A$
the settlement chain is carried by linear \code{Fulfill(ptr, x)} facts, each
consumed by exactly one downstream hop, so no party can be skipped. The
abstraction that buys unbounded path length also removes the collusion side
channel. Attacks that exist at the message layer must therefore be established
in $C$; only $A$ carries them to arbitrary $N$. This is a limitation of the
layering, and we state it rather than leaving it implicit.

\medbfparagraph{Tractability boundary}
One lemma, \code{Settle\_Requires\_Receiver\_Release}, remains open at the
abstract layer. Resolving \code{SenderSettled} backwards walks the entire
generic settle chain, whose length is unbounded, and the search does not
terminate. It verifies without difficulty in the fixed-hop theory. We report it
as the single genuine tractability boundary of this layer rather than removing
it from the file.
